\documentclass[12pt]{article}

% ---------------- PACKAGES ----------------
\usepackage{graphicx}
\usepackage[export]{adjustbox}
\usepackage{geometry}
\geometry{a4paper, margin=1in}
\usepackage{xcolor}
\usepackage{amsmath, amssymb}
\usepackage{booktabs}
\usepackage{array}
\usepackage{hyperref}
\usepackage[table]{xcolor}

\hypersetup{
    colorlinks=true,
    linkcolor=blue!60!black,
    urlcolor=blue!60!black
}

\usepackage{tikz}
\usetikzlibrary{arrows.meta, calc, positioning}
\usepackage{pgfplots}
\pgfplotsset{compat=1.18}

% ---------------- COLORS ----------------
\definecolor{deepblue}{RGB}{20,40,110}

% ---------------- FONTS ----------------
\usepackage{fontspec}
\setmainfont{TeX Gyre Termes}
\newfontfamily\titlefont{Libre Baskerville}

% ---------------- HEADER / FOOTER ----------------
\usepackage{fancyhdr}
\setlength{\headheight}{18pt}
\pagestyle{fancy}
\fancyhf{}
\fancyhead[L]{\small\color{deepblue} Olympiad Handout}
\fancyhead[R]{\small\color{deepblue} Kepler's Equation}
\fancyfoot[C]{\small Rajarshi Das Gupta \;|\; \thepage}
\renewcommand{\headrulewidth}{0.6pt}
\renewcommand{\footrulewidth}{0.6pt}

% ---------------- SPACING ----------------
\usepackage{setspace}
\setstretch{1.45}

% ---------------- SECTION STYLE ----------------
\usepackage{titlesec}
\titleformat{\section}{\large\bfseries\color{deepblue}}{{\thesection}}{0.8em}{}%
  [\vspace{-0.3em}\rule{\linewidth}{0.5pt}]
\titleformat{\subsection}{\normalsize\bfseries\color{deepblue}}
  {\thesubsection}{0.8em}{}

% ---------------- BOXES ----------------
\usepackage{tcolorbox}
\tcbuselibrary{skins, breakable}

\newenvironment{keyeq}{%
  \begin{tcolorbox}[
    enhanced, colback=blue!4!white, colframe=deepblue,
    boxrule=0.8pt, arc=4pt
  ]\centering
}{\end{tcolorbox}}

\newenvironment{notebox}[1]{%
  \begin{tcolorbox}[
    enhanced, breakable,
    colback=blue!3!white, colframe=deepblue,
    boxrule=0.7pt, arc=4pt,
    title={\bfseries #1}
  ]
}{\end{tcolorbox}}

\newenvironment{probbox}[2]{%
  \begin{tcolorbox}[
    enhanced, breakable,
    colback=blue!5!white, colframe=deepblue,
    boxrule=0.8pt, arc=4pt,
    title={\bfseries [#1] Problem #2}
  ]
}{\end{tcolorbox}}


% license 

\usepackage[
    type={CC},
    modifier={by-nc-nd},
    version={4.0},
]{doclicense}

\usepackage{hyperref}

\hypersetup{
    pdftitle={Keplar's Equation},
    pdfauthor={Rajarshi Das Gupta},
    pdfkeywords={LaTeX, Note, Copyright},
    pdfdisplaydoctitle=true,
    pdfstartview=FitH
}

% ===============================================================
\begin{document}

% ================= PAGE 1: COVER =================
\begin{titlepage}
\begin{center}

\vspace*{2cm}

{\small\scshape\color{deepblue} OLYMPIAD HANDOUT}

\vspace{2cm}

{\titlefont\fontsize{42}{44}\selectfont\bfseries\color{deepblue} Kepler's Equation}

\vspace{0.8cm}

{\large The Unsolvable Equation of Celestial Mechanics}

\vspace{1cm}
\rule{0.35\textwidth}{0.5pt}

\vspace{2cm}

{\large \textbf{Rajarshi Das Gupta}}

\vspace{0.3cm}

\vspace{2.5cm}

% ---------- QUOTE ----------
\begin{minipage}{0.7\textwidth}
\centering
\itshape
``I much prefer the sharpest criticism of a single intelligent man to the thoughtless approval of the masses.''\\[0.3cm]
--- Johannes Kepler
\end{minipage}

\vfill

{\small \today}

\end{center}
\end{titlepage}


% ================= PAGE 2: CONTENTS =================
\newpage

\begin{center}
{\Large\bfseries\color{deepblue} Contents}
\end{center}

\vspace{0.8cm}

\begin{center}
\begin{tabular}{|p{0.75\textwidth}|c|}
\hline

\rowcolor{deepblue}
\color{white}\textbf{Topics} & \color{white}\textbf{Page} \\
\hline

\textbf{1. Introduction} & 1 \\
\hline

\textbf{2. Geometry of the Ellipse} & 2 \\
\quad 2.1 Definition and Parameters & 2 \\
\quad 2.2 Focus and Eccentricity & 3 \\
\hline

\textbf{3. Auxiliary Circle} & 4 \\
\quad 3.1 Construction & 4 \\
\quad 3.2 Relation to Ellipse & 5 \\
\hline

\textbf{4. Area Law and Derivation} & 6 \\
\quad 4.1 Kepler’s Second Law & 6 \\
\quad 4.2 Area Calculation & 7 \\
\hline

\textbf{5. Kepler's Equation} & 8 \\
\hline

\textbf{6. Solving the Equation} & 9 \\
\quad 6.1 Iterative Methods & 9 \\
\quad 6.2 Convergence Issues & 10 \\
\hline

\textbf{7. Numerical Methods} & 11 \\
\quad 7.1 Newton-Raphson Method & 11 \\
\quad 7.2 Series Expansion & 12 \\
\hline

\textbf{8. Applications in Astronomy} & 13 \\
\hline

\textbf{9. Practice Problems} & 14 \\
\hline

\end{tabular}
\end{center}
\newpage


% ================= PAGE 3+: YOUR DOCUMENT STARTS =================





% ================================================================
\section{Introduction}
% ================================================================

We all know about Kepler's laws from our physics textbooks. But apart from giving
these three laws, Johannes Kepler is also known for deriving a remarkable equation
for analysing the motion of heavenly bodies, now known as \textbf{Kepler's Equation}.

Kepler's laws established that planets orbit in ellipses with the Sun at one focus,
that the radius vector sweeps equal areas in equal times, and that the orbital period
relates to orbital size. Yet these laws left one question unanswered: given an orbit
and an elapsed time, \emph{where exactly is the planet?} The non-uniform speed along
the ellipse means no simple formula expresses position as a function of time.
Kepler's Equation is the mathematical bridge that resolves this.

This article derives the equation from first principles using only Kepler's Second
Law and the geometry of the ellipse. The central tool is the \emph{auxiliary circle},
which reduces area calculations inside an ellipse to the simpler geometry of a circle.
Once derived, we explore the equation's transcendental nature and study the principal
numerical methods used to solve it in practice.

% ================================================================
\section{Geometry of Elliptical Orbits}
% ================================================================

\subsection{The Ellipse and Its Parameters}

An ellipse is characterised by its \textbf{semi-major axis} $a$ (the longer half-axis)
and its \textbf{semi-minor axis} $b$ (the shorter half-axis), related to the
\textbf{eccentricity} $e \in [0,1)$ by
\[
  b^2 = a^2(1 - e^2).
\]
When $e = 0$ the ellipse is a circle; as $e \to 1$ it becomes increasingly elongated.
The Earth has orbital eccentricity $e \approx 0.017$; Halley's Comet has orbital
eccentricity $e \approx 0.967$.

The two \textbf{foci} lie on the major axis at distance $ae$ from the centre $C$.
The Sun occupies one focus $F$. The closest orbital point to $F$ is the
\textbf{periapsis} at distance $a(1-e)$; the farthest is the \textbf{apoapsis}
at distance $a(1+e)$.

% ================================================================
\subsection{The Three Anomalies}
% ================================================================

To specify where a body is on its orbit, astronomers define angular quantities called
\textbf{anomalies}. The word comes from the Greek \emph{anomalia} (irregularity),
reflecting that planetary positions do not advance uniformly. Three anomalies are
used, shown together in Figure~1.

\vspace{0.4cm}
% ---- FIGURE 1 : signature diagram ----------------------------
\begin{center}
\begin{tikzpicture}[>=Stealth, line cap=round, line join=round, scale=1.15]
  % a=4, b=2.5, e=0.7, c=2.8
  % F at origin; C at (-2.8,0)
  % E=65 deg: P'=(-1.109, 3.625), P=(-1.109, 2.266)
  \coordinate (F)   at (0,      0);
  \coordinate (C)   at (-2.8,   0);
  \coordinate (Per) at (1.2,    0);
  \coordinate (Apo) at (-6.8,   0);
  \coordinate (Pp)  at (-1.109, 3.625);
  \coordinate (P)   at (-1.109, 2.266);
  \coordinate (Px)  at (-1.109, 0);

  \draw[->, gray!50] (-8.0,0)    -- (2.6,0)    node[right,font=\small]{$x$};
  \draw[->, gray!50] (-2.8,-3.0) -- (-2.8,3.5) node[above,font=\small]{$y$};

  \draw[thick, blue]     (C) ellipse [x radius=4, y radius=2.5];
  \draw[dashed, gray!60] (C) circle  [radius=4];

  \draw[dashed, gray!45] (Pp) -- (Px);
  \draw[thick, blue!55]  (C)  -- (Pp);
  \draw[thick, red!80]   (F)  -- (P);

  \fill[orange]  (F)  circle (3pt)   node[below right,font=\small]{$F$};
  \fill[black]   (C)  circle (2pt)   node[below left, font=\small]{$C$};
  \fill[blue!65] (Pp) circle (2.5pt) node[above right,font=\small]{$P'$};
  \fill[red!80]  (P)  circle (2.5pt) node[right,      font=\small]{$P$};

  \node[below right,font=\small] at (Per) {Periapsis};
  \node[below left, font=\small] at (Apo) {Apoapsis};

  \draw[|<->|] (C) -- (Per)            node[midway,fill=white,inner sep=1pt]{$a$};
  \draw[|<->|] (C) -- (-2.8, 2.5)      node[midway,left]{$b$};
  \draw[|<->|] (-2.8,-0.65)--(0,-0.65) node[midway,above]{$ae$};

  \draw[->, red!80, thick] (0.8,0)
      arc [start angle=0, end angle=120, radius=0.8];
  \node[red!80,font=\small] at (32:1.1) {$\nu$};

  \draw[->, blue!55, thick] (-2.8+1.0,0)
      arc [start angle=0, end angle=65, radius=1.0];
  \node[blue!55,font=\small] at (-2.8+0.62,0.55) {$E$};

  \draw[<->, gray!55, very thin]
      ($(P)+(0.22,0)$) -- ($(Pp)+(0.22,0)$)
      node[midway,right,font=\tiny,gray!55]{same $x$};
\end{tikzpicture}

\vspace{0.2cm}
{\small\textit{Figure~1: Elliptical orbit showing $\nu$ (red, at $F$),
$E$ (blue, at $C$), and the auxiliary circle (dashed). \\
Note: You can find an interective version of this diagram at  \href{https://rajarshineel.github.io/Anomalies-in-Ellipse/}{https://rajarshineel.github.io/Anomalies-in-Ellipse/}}}
\end{center}

\vspace{0.4cm}

\noindent
\textbf{True Anomaly $\nu$} is the angle at the focus $F$ from the periapsis direction
to the body $P$, measured in the direction of motion. It is the physically observable
position angle but grows \emph{non-uniformly} in time: fast near periapsis, slow near
apoapsis. No closed formula gives $\nu$ as a function of $t$ directly.

\bigskip\noindent
\textbf{Eccentric Anomaly $E$} is defined via the auxiliary circle (radius $a$,
centred at $C$). The body at $P$ is projected vertically onto the auxiliary circle
to give $P'$; the central angle $\angle P'CA$ is $E$. It is the key geometric
intermediary: any area in the ellipse equals $b/a$ times the corresponding area in
the auxiliary circle.
\newpage
\bigskip\noindent
\textbf{Mean Anomaly $M$} is defined as $M = 2\pi t/T$. It grows \emph{uniformly}
with time and is trivially computable from $t$, but does not directly locate the body
on the ellipse. Kepler's Equation converts $M$ into $E$, from which $\nu$ follows.

\vspace{0.5cm}
\begin{center}
\renewcommand{\arraystretch}{1.5}
\begin{tabular}{llll}
  \toprule
  \textbf{Anomaly} & \textbf{Symbol} & \textbf{Measured at} & \textbf{Nature} \\
  \midrule
  Mean Anomaly      & $M$ & Fictitious (from time)   & Uniform in time \\
  Eccentric Anomaly & $E$ & Centre $C$, aux.\ circle & Geometric intermediary \\
  True Anomaly      & $\nu$ & Focus $F$              & Physical (non-uniform) \\
  \bottomrule
\end{tabular}
\end{center}
\vspace{0.5cm}
% ================================================================
\section{Derivation of Kepler's Equation}
% ================================================================

\subsection{Kepler's Second Law and the Area--Time Relation}

The derivation rests on Kepler's Second Law: \\

\begin{notebox}{Kepler's Second Law}
{\textit{The radius vector from the focus to the planet sweeps out equal areas
in equal intervals of time.}}
\end{notebox}
\vspace{0.2cm}
\noindent

This follows from conservation of angular momentum in any central-force field.
Its practical consequence is that tracking area is equivalent to tracking time.

Let $T$ be the orbital period and $t$ the elapsed time since periapsis. Since the
total ellipse area is $\pi ab$ and swept area is proportional to time:
\[
  \frac{A}{\pi ab} = \frac{t}{T}
  \qquad\Longrightarrow\qquad
  A = \pi ab \cdot \frac{t}{T}.
\]
Figure~2 illustrates the swept area $A$.
\newpage

\vspace{0.3cm}
\begin{center}
\begin{tikzpicture}[scale=1.15]
  \coordinate (C)  at (0,      0);
  \coordinate (F)  at (2.4,    0);
  \coordinate (P0) at (4.0,    0);
  \coordinate (P)  at (-1.368, 2.349);

  \fill[orange!35]
    (F) -- (P0)
    arc [start angle=0, end angle=110, x radius=4, y radius=2.5]
    -- cycle;

  \draw[thick, blue] (C) ellipse [x radius=4, y radius=2.5];
  \draw[thick, black!80] (F) -- (P0);
  \draw[thick, black!80] (F) -- (P);

  \fill[orange]  (F)  circle (2.5pt) node[below right,font=\small]{$F$};
  \fill[black]   (P0) circle (2.0pt) node[below right,font=\small]{$P_0$ (periapsis)};
  \fill[blue!80] (P)  circle (2.0pt) node[above left, font=\small]{$P$ (body)};
  \node[font=\normalsize] at (1.6, 1.3) {$A$};
\end{tikzpicture}

\vspace{0.2cm}
{\small\textit{Figure~2: Swept area $A$ from periapsis.}}
\end{center}

\subsection{Introducing the Mean Anomaly}

Replace $t$ with $M = 2\pi t/T$, growing at constant rate from $0$ to $2\pi$.
The area--time relation becomes:
\[
  A = \frac{ab}{2}\,M. \tag{1}
\]
This is the first expression for $A$, linear in time. We now derive a second,
geometric expression for $A$ in terms of $E$, then equate the two.

\subsection{The Auxiliary Circle and Area Scaling}

The auxiliary circle (radius $a$, centred at $C$) is related to the ellipse by a
uniform vertical compression. Any circle point $(x,\,y_{\text{circle}})$ maps to
the ellipse point $(x,\,y_{\text{ellipse}})$ with
\[
  y_{\text{ellipse}} = \frac{b}{a}\,y_{\text{circle}}.
\]
Since every vertical strip scales by $b/a$, every ellipse area equals the
corresponding circle area scaled by $b/a$:
\[
  A_{\text{ellipse}} = \frac{b}{a}\,A_{\text{circle}}. \tag{2}
\]

\newpage
\subsection{Area in the Auxiliary Circle: Sector Minus Triangle}

The area swept in the auxiliary circle from periapsis $A$ to the projected point
$P'$ is sector $CAP'$ minus triangle $CFP'$ (Figure~3).
\vspace{0.4cm}

\begin{center}
\begin{tikzpicture}[scale=1.35]
  \coordinate (C)  at (0,      0);
  \coordinate (F)  at (-1.664, 0);
  \coordinate (A)  at (3.2,    0);
  \coordinate (Pp) at (1.835,  2.621);
  \coordinate (P)  at (1.835,  2.197);
  \coordinate (N)  at (1.835,  0);

  \fill[blue!55]
    (C) -- (A)
    arc [start angle=0, end angle=55, radius=3.2]
    -- cycle;
  \fill[red!55] (C) -- (F) -- (Pp) -- cycle;

  \draw[thick, gray!55, dashed] (C) circle [radius=3.2];
  \draw[thick, blue!80] (C) ellipse [x radius=3.2, y radius=2.681];

  \draw[thick, black!70] (C)  -- (A);
  \draw[thick, black!70] (C)  -- (Pp);
  \draw[thick, red!80]   (F)  -- (Pp);
  \draw[dashed, gray!60] (Pp) -- (N);
  \draw[dashed, gray!60] (P)  -- (N);

  \fill[black]     (C)  circle (2.5pt) node[below left, font=\small]{$C$};
  \fill[orange!90] (F)  circle (2.5pt) node[below,      font=\small]{$F$};
  \fill[black]     (Pp) circle (2.5pt) node[above right,font=\small]{$P'$};
  \fill[blue!80]   (P)  circle (2.5pt) node[right,      font=\small]{$P$};
  \fill[black]     (A)  circle (2pt)   node[below right,font=\small]{$A$};
  \fill[gray!70]   (N)  circle (1.8pt) node[below,      font=\small]{$N$};

  \draw[->, thick, black!70] (0.95,0)
      arc [start angle=0, end angle=55, radius=0.95];
  \node[font=\small, black] at (28:1.25) {$E$};

  \draw[|<->|,thick, black!70] (0,-0.58) -- (-1.664,-0.58)
      node[midway, below, font=\small, black!70]{$ae$};
  \draw[|<->|,thick, red!65] ($(Pp)+(0.42,0)$) -- ($(N)+(0.42,0)$)
      node[midway, left , font=\small,thick, red!65]{$a\sin E$};

\fill[red!55]   (4.1, 1.85) rectangle (4.55, 2.30);
  \draw[gray!60]  (4.1, 1.85) rectangle (4.55, 2.30);
  \node[right, font=\small, red!70!black] at (4.62, 2.07)
      {Triangle $CFP' = \tfrac{1}{2}a^2 e\sin E$};

  \fill[blue!55]  (4.1, 2.55) rectangle (4.55, 3.00);
  \draw[gray!60]  (4.1, 2.55) rectangle (4.55, 3.00);
  \node[right, font=\small, blue!70!black] at (4.62, 2.77)
      {Sector $CAP' = \tfrac{1}{2}a^2 E$};

  

  \node[font=\small, align=left, black] at (5.30, 1.20)
      {Swept area $= \tfrac{1}{2}a^2(E - e\sin E)$};
\end{tikzpicture}

\vspace{0.2cm}
{\small\textit{Figure~3: Area decomposition: sector (blue) minus triangle (pink).}}
\end{center}

\vspace{0.3cm}
\noindent
\textbf{Sector $CAP'$:} radius $a$, angle $E$, area $= \tfrac{1}{2}a^2 E$.

\noindent
\textbf{Triangle $CFP'$:} base $|CF| = ae$, height $P'N = a\sin E$,
area $= \tfrac{1}{2}a^2 e\sin E$.

\noindent
\textbf{Why sector minus triangle?}  The swept area is measured from the focus $F$,
not the centre $C$. Since $F$ is displaced from $C$ by $ae$, the triangle $CFP'$
lies inside the sector but outside the region swept from $F$; subtracting it gives
the correct area.

\noindent The swept area in the auxiliary circle is:
\[
  A_{\text{circle}} = \tfrac{1}{2}a^2 E - \tfrac{1}{2}a^2 e\sin E
                    = \frac{a^2}{2}(E - e\sin E).
\]

\subsection{Scaling Back to the Ellipse}

Applying relation (2):
\[
  A = \frac{b}{a} \cdot \frac{a^2}{2}(E - e\sin E)
    = \frac{ab}{2}(E - e\sin E). \tag{3}
\]
\newpage
\subsection{Equating Both Expressions}

Expressions (1) and (3) both equal the swept area $A$. Setting them equal and
cancelling $\tfrac{ab}{2}$:

\vspace{-0.4cm}
\begin{keyeq}
\boxed{M = E - e\sin E}
\end{keyeq}
\vspace{0.2cm}
\noindent

This is \textbf{Kepler's Equation}.\hfill$\blacksquare$

\vspace{0.4cm}
\begin{notebox}{Important: Angles Must Be in Radians}
The sector area formula $\tfrac{1}{2}a^2 E$ is valid only when $E$ is in
\textbf{radians}. Both $E$ and $M$ must be in radians; using degrees introduces
a spurious factor of $\pi/180$.
\end{notebox}

% ================================================================
\section{From Eccentric to True Anomaly}
% ================================================================

Kepler's Equation yields $E$ from $M$. We still need to convert $E$ into the
physically meaningful true anomaly $\nu$ and the radial distance $r$.

\subsection{Coordinate Relations}

A point $P$ on the ellipse has parametric coordinates from centre $C$:
\[
  x = a\cos E, \qquad y = b\sin E.
\]
In polar form $(r,\nu)$ from focus $F$ (offset from $C$ by $ae$):
\[
  x = r\cos\nu + ae, \qquad y = r\sin\nu.
\]

\subsection{Radial Distance}

Equating both expressions and eliminating $b$ via $b^2 = a^2(1-e^2)$:
\[
  \boxed{r = a\,(1 - e\cos E)}.
\]
Equivalently in terms of $\nu$:
\[
  r = \frac{a(1-e^2)}{1 + e\cos\nu}.
\]

\subsection{Half-Angle Conversion}

From the $x$-expressions with $r = a(1-e\cos E)$:
$\cos E = (e + \cos\nu)/(1 + e\cos\nu)$.
Applying the half-angle tangent identity:
\[
  \boxed{\tan\frac{E}{2} \;=\; \sqrt{\frac{1-e}{1+e}}\;\tan\frac{\nu}{2}}.
\]
This is strictly monotone on $[0, 2\pi]$: each $E$ gives a unique $\nu$.

\vspace{0.4cm}
\begin{center}
\renewcommand{\arraystretch}{1.6}
\begin{tabular}{llll}
  \toprule
  \textbf{Anomaly} & $x$ \textbf{from} $C$ & $x$ \textbf{from} $F$ & $y$ \\
  \midrule
  True ($\nu$)    & $r\cos\nu$  & $r\cos\nu - ae$  & $r\sin\nu$ \\
  Eccentric ($E$) & $a\cos E$   & $a(\cos E - e)$  & $b\sin E$  \\
  \bottomrule
\end{tabular}
\end{center}

% ================================================================
\section{Understanding and Visualising Kepler's Equation}
% ================================================================

\subsection{What Does Kepler's Equation Actually Tell Us?}

Kepler's Equation $M = E - e\sin E$ cannot be solved algebraically for $E$ given
$M$. This is because $E$ appears both inside and outside the sine function --- the
equation is \textbf{transcendental}. Kepler himself recognised this and spent years
searching for a closed-form solution, eventually concluding that none existed.

Since no formula gives $E$ directly from $M$, we must rely on \textbf{numerical
methods}: iterative procedures that start with an approximate value and improve it
step by step until the result is accurate to any desired precision.

The solution chain from time to position is therefore:
\[
  t
  \;\xrightarrow{M = 2\pi t/T}\;
  M
  \;\xrightarrow{M = E - e\sin E}\;
  E
  \;\xrightarrow{\tan(E/2) = \cdots}\;
  \nu
  \;\xrightarrow{r = a(1-e\cos E)}\;
  (r,\,\nu).
\]
Steps 1, 3, and 4 are direct arithmetic. Only Step~2 requires iteration.

\subsection{The Geometric Meaning: $M$, $E$, and $\nu$ as Angles on the Same Orbit}

Figure~4 shows how $M$, $E$, and $\nu$ relate to each other geometrically for
a single orbit position. The key insight is:

\begin{itemize}
  \item $M$ is where a \emph{fictitious} body would be, moving at constant speed
        on the auxiliary circle.
  \item $E$ is the angle to the projection $P'$ of the real body onto the
        auxiliary circle, measured at the centre $C$.
  \item $\nu$ is the angle to the real body $P$, measured at the focus $F$.
\end{itemize}

For a circular orbit ($e = 0$) all three are identical at all times. For an
elliptical orbit, they diverge: near periapsis the planet moves faster than the
fictitious body ($\nu > M$), and near apoapsis it moves slower ($\nu < M$).
Kepler's Equation quantifies exactly this departure.

\vspace{0.4cm}
\begin{center}
\begin{tikzpicture}[scale=1.1]
  \coordinate (C)   at (0,      0);
  \coordinate (F)   at (-1.925, 0);
  \coordinate (Per) at (3.5,    0);
  \coordinate (Pp)  at (2.250,  2.681);
  \coordinate (P)   at (2.250,  2.247);
  \coordinate (Mb)  at (3.150,  1.525);

  \draw[->, gray!50] (-5.6,0)   -- (4.2,0)   node[right,font=\small]{$x$};
  \draw[->, gray!50] (-1.925,-3.2)--(-1.925,3.5) node[above,font=\small]{};

  \draw[thick, blue]     (C) ellipse [x radius=3.5, y radius=2.934];
  \draw[dashed, gray!60] (C) circle  [radius=3.5];

  \draw[thick, green!50!black] (C) -- (Mb);
  \draw[thick, blue!55]        (C) -- (Pp);
  \draw[thick, red!80]         (F) -- (P);
  \draw[dashed, gray!40]       (Pp) -- (2.250,0);

  \fill[orange]          (F)   circle (2.5pt) node[below,       font=\small]{$F$};
  \fill[black]           (C)   circle (2pt)   node[below left,  font=\small]{$C$};
  \fill[green!50!black]  (Mb)  circle (2.5pt) node[right,       font=\small]{$M$-body};
  \fill[blue!65]         (Pp)  circle (2.5pt) node[above right, font=\small]{$P'$};
  \fill[red!80]          (P)   circle (2.5pt) node[right,       font=\small]{$P$};
  \fill[black]           (Per) circle (1.8pt) node[below right, font=\small]{Periapsis};

  \draw[->, green!50!black, thick] (0.8,0)
      arc [start angle=0, end angle=25.86, radius=0.8];
  \node[green!50!black, font=\small] at (11:1.4) {$M$};

  \draw[->, blue!55, thick] (1.05,0)
      arc [start angle=0, end angle=50, radius=1.05];
  \node[blue!55, font=\small] at (35:1.4) {$E$};

  \draw[->, red!80, thick]
      ($( F)+(1.0,0)$) arc [start angle=0, end angle=29.5, radius=1.0];
  \node[red!80, font=\small] at ($( F)+(18:1.4)$) {$\nu$};

  \node[font=\small, align=center] at (0,-4)
      {Near periapsis: $\nu > E > M$};
\end{tikzpicture}

\vspace{0.2cm}
{\small\textit{Figure~4: Geometric comparison of $M$ (green), $E$ (blue), and $\nu$
(red) at the same orbital position near periapsis. The real body $P$ is ahead of
both the $E$-projection and the mean-motion body.}}
\end{center}
\newpage 
% ================================================================
\section{Numerical Methods for Solving Kepler's Equation}
% ================================================================

\subsection{Why No Closed Form Exists}

The equation $M = E - e\sin E$ is called \textbf{transcendental}: $E$ appears both
as a plain number and hidden inside $\sin E$ at the same time. No matter how cleverly
you rearrange it, you cannot write $E$ directly as a formula in $M$ and $e$. Kepler
himself spent years trying and gave up. This is not a gap in our knowledge --- it is
a proven mathematical impossibility.

What we can do instead is \emph{guess} a value of $E$, check how far off it is, and
correct it. We repeat this until the answer is as accurate as we need. These
step-by-step correction procedures are called \textbf{numerical methods}\footnote{For a deep dive in numerical methods, watch this playlist by Jeffrey Chasnov. (\hyperlink{https://youtube.com/playlist?list=PLkZjai-2Jcxn35XnijUtqqEg0Wi5Sn8ab&si=vtjDXGNt5P7oDNnd}{link})}. The good news is that Kepler's equation always has \emph{exactly one} answer (because the
left-hand side grows faster than the right-hand side for any $e < 1$), so every
reasonable method will eventually find it. 

\subsection{Method 1: Fixed-Point Iteration}

The idea is beautifully simple. Rewrite Kepler's equation by moving $e\sin E$ to the
right:
\[
  E = M + e\sin E.
\]
This says: ``the answer $E$ equals $M$ plus a small correction $e\sin E$.''  We use
this as a recipe. Start with the guess $E_0 = M$ (pretend the correction is zero),
then keep plugging the current value back in:

\begin{keyeq}
  $E_{n+1} = M + e\sin E_n$
\end{keyeq}

Each step, the correction gets better. This works because the correction term
$e\sin E$ is always smaller than $e < 1$ times the current error, so the error
shrinks by a factor of roughly $e$ every step. For planets with small eccentricity
this is very fast; for highly elongated orbits it takes more steps.
\newpage 

\begin{notebox}{Quick Example --- Fixed-Point}
Earth: $e = 0.0167$, $M = 1.0\,\text{rad}$.
\begin{align*}
  E_0 &= 1.0000 \\
  E_1 &= 1.0 + 0.0167 \times \sin(1.0) = 1.0 + 0.0167 \times 0.8415 = 1.01405 \\
  E_2 &= 1.0 + 0.0167 \times \sin(1.01405) \approx 1.01406
\end{align*}
Two steps give six-digit accuracy. The correction is tiny because Earth's orbit is
nearly circular.
\end{notebox}

\subsection{Method 2: Newton--Raphson Method}

Fixed-point iteration is fine for nearly circular orbits, but it slows down when $e$
is large. The Newton--Raphson method is smarter: instead of just applying the
correction blindly, it looks at \emph{how steeply} the equation is changing and
adjusts the step size accordingly. The idea is to draw a tangent line to the error
curve and find where that line crosses zero --- that crossing is a much better next
guess.

Define the ``error'' at a given guess as $f(E) = E - e\sin E - M$. We want $f(E) = 0$.
The Newton--Raphson update is:

\begin{keyeq}
  $\displaystyle E_{n+1} = E_n - \frac{E_n - e\sin E_n - M}{1 - e\cos E_n}$
\end{keyeq}

The numerator is the current error; the denominator is the slope of $f$ at $E_n$.
Dividing by the slope scales the correction to the right size. The result is that
the number of accurate decimal digits \emph{doubles} with every step --- this is
called quadratic convergence, and it is far faster than fixed-point iteration.

\begin{notebox}{Quick Example --- Newton--Raphson}
Same case: $e = 0.0167$, $M = 1.0\,\text{rad}$, start with $E_0 = 1.0$.
\[
  E_1 = 1.0 - \frac{1.0 - 0.0167\sin(1.0) - 1.0}{1 - 0.0167\cos(1.0)}
      = 1.0 - \frac{-0.01405}{0.9991} \approx 1.01406.
\]
One single step reaches the same accuracy that fixed-point needed two steps for ---
and for a high-eccentricity orbit the gap would be far more dramatic.
\end{notebox}

\begin{center}
\renewcommand{\arraystretch}{1.5}
\begin{tabular}{lll}
  \toprule
  \textbf{Method} & \textbf{Speed} & \textbf{Best used for} \\
  \midrule
  Fixed-point    & Error $\times\,e$ per step (linear)   & Low-eccentricity orbits ($e < 0.3$) \\
  Newton--Raphson & Digits double per step (quadratic)    & All orbits; standard choice \\
  \bottomrule
\end{tabular}
\end{center}

\begin{notebox}{Practical Guidelines}
\begin{itemize}
  \item Both methods require $E$ and $M$ in \textbf{radians}, not degrees.
  \item A good starting guess is always $E_0 = M$. For very eccentric orbits
        ($e > 0.8$), starting with $E_0 = \pi$ avoids occasional slow convergence.
  \item Newton--Raphson is the standard choice in real astronomical software. It
        converges to full computer precision in 3--5 iterations for any physically
        meaningful orbit.
\end{itemize}
\end{notebox}

\subsection{A Short Worked Example: Finding a Planet's Position}

An exoplanet has eccentricity $e = 0.30$ and period $T = 2.0$ years.
Find its eccentric anomaly $0.5$ years after periapsis.

\textbf{Step 1 --- Mean anomaly:}
\[
  M = \frac{2\pi}{T} \times t = \frac{2\pi}{2.0} \times 0.5 = \frac{\pi}{2}
    \approx 1.5708\,\text{rad}.
\]

\textbf{Step 2 --- Newton--Raphson} (two steps from $E_0 = M = 1.5708$):
\[
  E_1 = 1.5708 - \frac{1.5708 - 0.30\sin(1.5708) - 1.5708}{1 - 0.30\cos(1.5708)}
      = 1.5708 - \frac{-0.3000}{1.000} = 1.8708\,\text{rad}.
\]
\[
  E_2 = 1.8708 - \frac{1.8708 - 0.30\sin(1.8708) - 1.5708}{1 - 0.30\cos(1.8708)}
      = 1.8708 - \frac{0.0263}{1.0592} \approx 1.8460\,\text{rad}.
\]
\[
  E_3 \approx 1.8453\,\text{rad} \quad\text{(converged to four decimal places)}.
\]

\textbf{Step 3 --- True anomaly:}
\[
  \tan\frac{\nu}{2} = \sqrt{\frac{1+0.30}{1-0.30}}\,\tan\frac{1.8453}{2}
                    = \sqrt{1.857} \times \tan(0.9227) \approx 1.363 \times 1.297
                    \approx 1.768,
\]
\[
  \frac{\nu}{2} \approx 60.5^\circ \quad\Rightarrow\quad \nu \approx 121^\circ.
\]

\textbf{Step 4 --- Heliocentric distance} (taking $a = 1$ for simplicity):
\[
  r = a(1 - e\cos E) = 1 - 0.30\cos(1.8453) \approx 1 - 0.30\times(-0.2756)
    \approx 1.083\,a.
\]
The planet is about $8\%$ farther from its star than the semi-major axis, which makes
sense: at a true anomaly of $121^\circ$ it has moved well past the halfway point of
its orbit and is receding toward apoapsis.


% ================================================================
%  PROBLEMS
% ================================================================

\section*{Problems}
\addcontentsline{toc}{section}{Problems}

% ------------------------------------------------------------------
\begin{probbox}{5 points}{1. A Simple Circular Check}
A planet moves in a perfectly circular orbit with eccentricity $e = 0$ and period
$T = 1\,\text{year}$.
\begin{enumerate}
  \item[(a)] Find the mean anomaly $M$ exactly 3 months after periapsis.
             Give your answer in both radians and degrees.
  \item[(b)] State the eccentric anomaly $E$ without any iteration.
             Explain why no iteration is needed.
  \item[(c)] What is the true anomaly $\nu$, and where is the planet on its orbit?
\end{enumerate}

\medskip
\textbf{\color{blue!60!black}Answer.}

\textbf{(a)} Three months is one quarter of the period, so
\begin{equation*}
  M \;=\; \frac{2\pi}{T}\,t \;=\; \frac{2\pi}{1}\times\frac{1}{4}
        \;=\; \frac{\pi}{2} \;\approx\; 1.571\,\text{rad} \;=\; 90^\circ.
\end{equation*}

\textbf{(b)} Setting $e = 0$ in Kepler's equation gives $M = E - 0\cdot\sin E = E$,
so $E = \pi/2$ exactly. No iteration is needed because for a circular orbit the
equation is already solved --- $M$, $E$, and $\nu$ are all identical at every point.

\textbf{(c)} The half-angle formula with $e = 0$ gives
$\tan(\nu/2) = \tan(E/2)$, hence $\nu = E = 90^\circ$.
The planet is exactly one quarter of the way around its orbit, at the tip of the
semi-minor axis, equidistant from periapsis and apoapsis.
\end{probbox}

\vspace{0.6cm}

% ------------------------------------------------------------------
\begin{probbox}{10 points}{2. Mars at Opposition}
Mars has eccentricity $e = 0.093$, period $T = 686.97\,\text{days}$, and
semi-major axis $a = 1.524\,\text{au}$. Assume Mars was at perihelion on
1 January 2023. The Mars opposition of 13 October 2023 occurred 285 days after
perihelion.
\begin{enumerate}
  \item[(a)] Find the mean anomaly $M$ of Mars on 13 October 2023.
  \item[(b)] Use two Newton--Raphson iterations (starting from $E_0 = M$) to find
             the eccentric anomaly $E$ to four decimal places.
  \item[(c)] Convert $E$ to true anomaly $\nu$.
  \item[(d)] Find Mars's distance from the Sun on that date.
  \item[(e)] Is the October 2023 opposition favourable for observation? Justify
             briefly by comparing $r$ with the perihelion and aphelion distances.
\end{enumerate}

\medskip
\textbf{\color{blue!60!black}Answer.}

\textbf{(a)}
\begin{equation*}
  M \;=\; \frac{2\pi \times 285}{686.97} \;\approx\; 2.6079\,\text{rad}
        \;\approx\; 149.4^\circ.
\end{equation*}

\textbf{(b)} Newton--Raphson with $E_0 = 2.6079$:
\begin{align*}
  E_1 &= 2.6079 - \frac{2.6079 - 0.093\sin(2.6079) - 2.6079}{1 - 0.093\cos(2.6079)}
       = 2.6079 - \frac{-0.0465}{1.0793} \approx 2.6510,\\[4pt]
  E_2 &= 2.6510 - \frac{2.6510 - 0.093\sin(2.6510) - 2.6079}{1 - 0.093\cos(2.6510)}
       \approx 2.6510 - \frac{0.0003}{1.0804} \approx 2.6507.
\end{align*}
Converged value: $E \approx 2.6507\,\text{rad}$.

\textbf{(c)} Using the half-angle formula:
\begin{equation*}
  \tan\frac{\nu}{2}
    = \sqrt{\frac{1+0.093}{1-0.093}}\,\tan\!\left(\frac{2.6507}{2}\right)
    = \sqrt{1.205}\times\tan(1.3254)
    \approx 1.098\times 4.018
    \approx 4.412,
\end{equation*}
\begin{equation*}
  \nu \approx 2\arctan(4.412) \approx 154.3^\circ.
\end{equation*}

\textbf{(d)}
\begin{equation*}
  r = a\,(1 - e\cos E)
    = 1.524\times\bigl(1 - 0.093\times\cos(2.6507)\bigr) 
    \end{equation*}
\begin{equation*}
    = 1.524\times\bigl(1 + 0.093\times 0.877\bigr)
    \approx 1.648\,\text{au}.
\end{equation*}

\textbf{(e)} Perihelion distance: $a(1-e) = 1.524\times 0.907 = 1.382\,\text{au}$.
Aphelion distance: $a(1+e) = 1.524\times 1.093 = 1.666\,\text{au}$.
At $r \approx 1.648\,\text{au}$, Mars is very close to aphelion, making this
opposition \emph{unfavourable} --- Mars is nearly as far from Earth as possible.
The best oppositions occur when Mars is near perihelion (around August--September).
\end{probbox}

\vspace{0.6cm}

% ------------------------------------------------------------------
\begin{probbox}{15 points}{3. Comet Time of Flight}
A comet has semi-major axis $a = 10\,\text{au}$ and eccentricity $e = 0.80$.
It is observed at true anomaly $\nu_1 = 60^\circ$, and observed again at
$\nu_2 = 240^\circ$.
Use $G = 6.67\times10^{-11}\,\text{N\,m}^2\,\text{kg}^{-2}$,
$M_\odot = 2.0\times10^{30}\,\text{kg}$, and
$1\,\text{au} = 1.496\times10^{11}\,\text{m}$.
\begin{enumerate}
  \item[(a)] Convert $\nu_1$ and $\nu_2$ to eccentric anomalies $E_1$ and $E_2$.
  \item[(b)] Find the mean anomalies $M_1$ and $M_2$.
  \item[(c)] Find the orbital period $T$.
  \item[(d)] Find the time elapsed between the two observations.
  \item[(e)] Find the comet's distance from the Sun at each observation.
             At which point is it moving faster, and why?
\end{enumerate}

\medskip
\textbf{\color{blue!60!black}Answer.}

\textbf{(a)} The conversion factor is
\begin{equation*}
  k = \sqrt{\frac{1-e}{1+e}} = \sqrt{\frac{0.20}{1.80}} \approx 0.3333.
\end{equation*}
For $\nu_1 = 60^\circ$:
\begin{equation*}
  \tan\frac{E_1}{2} = 0.3333\times\tan 30^\circ = 0.3333\times 0.5774 = 0.1925
  \;\Rightarrow\; E_1 \approx 21.8^\circ \approx 0.3802\,\text{rad}.
\end{equation*}
For $\nu_2 = 240^\circ$, we have $\tan 120^\circ = -\sqrt{3}$, and since
$\nu_2 > 180^\circ$ the comet is past apoapsis, placing $E_2$ in the range
$(180^\circ, 360^\circ)$:
\begin{equation*}
  \tan\frac{E_2}{2} = 0.3333\times(-\sqrt{3}) = -0.5774
  \;\Rightarrow\; E_2 \approx 300^\circ \approx 5.2360\,\text{rad}.
\end{equation*}

\textbf{(b)} Applying Kepler's equation $M = E - e\sin E$:
\begin{align*}
  M_1 &= 0.3802 - 0.80\times\sin(0.3802)
       = 0.3802 - 0.80\times 0.3714
       = 0.0831\,\text{rad},\\[4pt]
  M_2 &= 5.2360 - 0.80\times\sin(5.2360)
       = 5.2360 - 0.80\times(-0.8660)
       = 5.9288\,\text{rad}.
\end{align*}

\textbf{(c)} Converting the semi-major axis:
$a = 10\times 1.496\times10^{11} = 1.496\times10^{12}\,\text{m}$.
\begin{equation*}
  T = 2\pi\sqrt{\frac{a^3}{GM_\odot}}
    = 2\pi\sqrt{\frac{(1.496\times10^{12})^3}{6.67\times10^{-11}\times 2.0\times10^{30}}}
    \approx 2\pi \times 1.584\times10^{8}\,\text{s}
    \approx 31.6\,\text{years}.
\end{equation*}

\textbf{(d)}
\begin{equation*}
  \Delta t = \frac{M_2 - M_1}{2\pi}\times T
           = \frac{5.9288 - 0.0831}{6.2832}\times 31.6
           = 0.930 \times 31.6
           \approx 29.4\,\text{years}.
\end{equation*}

\textbf{(e)}
\begin{align*}
  r_1 &= a(1 - e\cos E_1)
       = 10\,(1 - 0.80\times\cos 21.8^\circ)
       = 10\,(1 - 0.80\times 0.929)
       = 2.57\,\text{au},\\[4pt]
  r_2 &= a(1 - e\cos E_2)
       = 10\,(1 - 0.80\times\cos 300^\circ)
       = 10\,(1 - 0.80\times 0.500)
       = 6.00\,\text{au}.
\end{align*}
The comet moves faster at the first observation ($r_1 = 2.57\,\text{au}$) because
it is closer to the Sun. By Kepler's second law, the radius vector sweeps equal
areas in equal times, so the speed must be higher when the distance is smaller.
\end{probbox}

\subsection*{For further exploration:}
1) A Textbook on Spherical Astronomy by W. M. Smart \href{https://archive.org/details/dli.ernet.240942/page/98/mode/2up}{(Chapter V)} \\
2) This YouTube video by Welch Labs: \href{https://www.youtube.com/watch?v=hBkmyJ3TE0g}{Kepler’s Impossible Equation}

\vfill
\doclicenseThis
\end{document}
